\section{Troubleshooting}\label{sec:troubleshooting}

\begin{description}[style=nextline, leftmargin=0pt, labelwidth=0pt, labelsep=0pt]
  \item[No trace or a stale trace] Check the source signal, cable, input
  gain, and Run light. Ensure LO Freq is below HI Freq. If you changed
  Length or loaded a stopped patch, resume capture to fill the buffer.
  \item[Peaks are too broad or move slowly] Reduce Smooth and Average.
  Compare window functions. Increasing Length narrows the bin spacing but
  also increases the time represented in each frame.
  \item[Levels differ from another analyzer] Match input gain, window,
  FFT length, smoothing, and amplitude reference. Set Slope to zero and
  account for summed polyphonic voices. A spectral peak is not a broadband
  RMS measurement.
  \item[Display reaches the top] Reduce the relevant input gain or the
  positive slope. The graph can exceed its visible range even though the
  source remains unchanged.
  \item[Analysis uses too much CPU] Try a longer Hop or shorter Length.
  Fourier performs analysis on Rack's audio engine thread despite having no
  audio outputs. Stopping Run stops capture but does not stop FFT processing.
\end{description}
